\section{Conclusiones}

El proyecto evaluó la posibilidad de anticipar cancelaciones de reservas
hoteleras y utilizar las probabilidades estimadas para priorizar acciones
preventivas. La separación temporal entre desarrollo (2015--2016) y holdout
(2017) permitió observar una dificultad que una partición aleatoria habría
ocultado parcialmente: todos los modelos presentaron una caída de rendimiento
al generalizar hacia un período posterior.

XGBoost obtuvo el mejor ordenamiento del target original, con ROC AUC 0,867291
sobre 2017 y Accuracy 0,781724. Logistic Regression mejorada alcanzó ROC AUC
0,855017, el mayor F1 (0,713461) y una estructura más interpretable. La MLP
consiguió el mayor Recall (0,964433) con threshold 0,5, pero a costa de 14.611
falsos positivos. El árbol resultó especialmente útil para estudiar y reducir
variables, mientras que Random Forest y KNN no superaron a los candidatos
principales en el holdout temporal.

La ingeniería de variables produjo una mejora clara en el modelo lineal. La
representación por bandas de \texttt{LeadTime}, las solicitudes especiales y
el historial relativo de cancelaciones permitieron mejorar el baseline sin
utilizar el holdout para seleccionar alternativas. También se descartaron
variables redundantes o con riesgo de leakage, priorizando una representación
más parsimoniosa y disponible al momento de la predicción.

La evaluación operativa reveló que el objetivo técnico y el objetivo económico
no están completamente alineados. Los modelos distinguen adecuadamente las
cancelaciones generales, pero asignan scores mayores principalmente a
cancelaciones no críticas. Dentro de las reservas canceladas, los ROC AUC para
distinguir No-Show y cancelaciones tardías fueron inferiores a 0,5. Por este
motivo, XGBoost fue el mejor modelo técnico, mientras que Logistic Regression
produjo la mejor priorización económica entre los candidatos analizados.

El valor bruto de una reserva no se interpretó como pérdida observada. El
análisis final incorporó margen de contribución, fracciones no recuperadas por
tipo de evento, efectividad de la intervención y costo por contacto. Ninguna
política resultó rentable bajo los escenarios conservador e intermedio. Los
ahorros fueron positivos principalmente en el escenario favorable extremo y
dependieron de supuestos que el dataset no permite verificar. Además, no fue
posible cuantificar descuentos innecesarios, fricción ni cancelaciones
inducidas por el contacto. En consecuencia, el sistema no demuestra por sí
solo una fuente de ahorro y no justifica una implementación operativa sin
validación adicional.

\subsection{Limitaciones}

Los datos proceden de dos hoteles portugueses y cubren solamente el período
comprendido entre julio de 2015 y agosto de 2017, por lo que los resultados no
se generalizan automáticamente a otros establecimientos o contextos. Algunas
variables categóricas presentan niveles de muy baja frecuencia, y las
diferencias temporales observadas sugieren posibles cambios en la distribución
de los datos.

El dataset tampoco contiene el margen real, los pagos efectivamente
recuperados, la posibilidad de revender cada habitación, el costo de contacto
ni la respuesta causal a una intervención. Los escenarios económicos permiten
analizar condiciones de viabilidad, pero no estimar un ahorro contable real.
Finalmente, \texttt{IsCriticalCancellation} se construyó retrospectivamente a
partir del estado final y de su fecha; sólo puede utilizarse como target de un
nuevo entrenamiento, nunca como predictor disponible al crear la reserva.

\subsection{Trabajo futuro}

La principal extensión consiste en repetir el ciclo de modelado utilizando
\texttt{IsCriticalCancellation} como objetivo. Esto permitiría optimizar
directamente la priorización de No-Show y cancelaciones tardías en lugar de
esperar que un modelo de cancelación general las ordene adecuadamente.

Antes de una aplicación real también sería necesario estimar los parámetros
económicos con información del hotel y diseñar una prueba controlada de la
política de contacto. Dicha prueba debería medir confirmaciones, pagos
garantizados, reventa del inventario, satisfacción del cliente y posibles
cancelaciones inducidas. Con esos datos podrían ajustarse capacidad y
threshold según beneficio neto, no solamente mediante métricas de
clasificación.

Como extensiones técnicas adicionales pueden evaluarse calibración de
probabilidades, monitoreo de drift temporal y nuevas búsquedas de
hiperparámetros. Estas mejoras sólo tendrían sentido después de alinear el
target con la decisión operativa y disponer de costos observados.
